\documentclass[11pt]{article}
\usepackage[margin=1in]{geometry}
\usepackage{amsmath,amssymb,amsthm,mathtools}
\usepackage[T1]{fontenc}
\usepackage{lmodern}
\usepackage{microtype}
\usepackage[colorlinks=true,linkcolor=blue,urlcolor=blue,citecolor=blue]{hyperref}
\newtheorem{theorem}{Theorem}[section]
\newtheorem{proposition}[theorem]{Proposition}
\newtheorem{lemma}[theorem]{Lemma}
\newtheorem{corollary}[theorem]{Corollary}
\theoremstyle{definition}\newtheorem{definition}[theorem]{Definition}
\theoremstyle{remark}\newtheorem{remark}[theorem]{Remark}
\newcommand{\SLE}{\mathrm{SLE}_6}
\newcommand{\HH}{\mathbb H}
\newcommand{\CC}{\mathbb C}
\newcommand{\RR}{\mathbb R}
\newcommand{\Law}{\operatorname{Law}}
\newcommand{\nuA}[4]{\nu^{#1}_{#2,#3,#4}}
\DeclareMathOperator{\Id}{Id}

\title{Local exploration laws and regular driving characteristics}
\hypersetup{pdftitle={Local exploration laws and regular driving characteristics},pdfauthor={Autonomous mathematical research}}
\author{Autonomous mathematical investigation, with independent critical checks}
\date{26 October 2026}
\begin{document}
\maketitle
\begin{abstract}
We investigate the corrected source-neighborhood locality axiom.
For continuous semimartingale Loewner drivers with absolutely continuous,
right-continuous drift, locality itself gives an initial quadratic-variation
rate. Two shrinking mirrored slits show that local boundedness of the
initial drift forces zero drift and variance six. Strong domain Markov
then identifies SLE$_6$, without assuming conformal covariance or Brownian
driving. The exact instability rates are inverse-linear and inverse-cubic
in the slit parameter. Separately, exact scaling proves that a fixed
nonconformal affine SLE$_6$ deformation cannot have the stated tame drift
structure in every bounded domain. These are conditional results: the
analytic hypotheses have not been derived from RSW or from the project
axioms, and the corrected classification and universality remain open.
\end{abstract}

\section{Current locality and scope}
Two colored marked domains that agree near their common source must have
identical initial exploration laws until exit from that neighborhood,
even if their remote targets differ. This is the user's corrected locality
axiom at commit \texttt{0e47018}. We fix one color orientation throughout.
The terminal-germ construction for the earlier fixed-target axiom is
excluded; its original proofs are retained in
\texttt{archive/original-locality-results.tex} and the detailed research record.

The first theorem below needs only a deletion consequence of corrected
locality. It does not claim that locality supplies regular driving drift.
For rough restart domains we state the prime-end neighborhood or admissible
test-deletion requirement explicitly. All driving clocks use
$\partial_tg=2/(g-W)$ and hence $\operatorname{hcap}(K_t)=2t$.
The standard conformal-map comparison for Loewner chains is derived below;
ordinary SLE facts used for affine examples are described in \cite{lawler,yuan}.
Continuous Loewner generation and infinite lifetime of affine SLE images
were proved in the companion research record \cite{record}.
% Research fragment for the corrected source-neighborhood locality assumption.
% Requires amsmath/amssymb/amsthm; all additional hypotheses are explicit.
\section{A near-source drift-boundedness criterion}
\label{sec:strong-local-loewner}

\paragraph{Status.}
This is a new conditional rigidity result.  It does not assume conformal
invariance, Brownian driving, or local-martingale driving.  Instead it shows
that a locally bounded \emph{initial driving drift}, together with an
absolutely continuous drift part and locality, already forces zero drift and
variance six.  No regularity of quadratic variation is assumed.  Strong DMP then identifies the whole driving law.  The
boundedness and characteristic hypotheses are not derived from RSW or from
the three project axioms.  A finite-dimensional bulk-jet closure is given as
one explicit sufficient hypothesis.

Only a particular deletion consequence of the corrected source-neighborhood
locality is needed.  Arbitrary target changes and the full force of corrected
locality are not used in this theorem.  Thus the result is compatible with,
but does not exhaust, the corrected research objective.

\subsection{Chart-dependent initial characteristics and their covariance}

Let $f:\mathbb H\to D$ be an inverse chordal conformal chart, with source
$f(0)=p$ and target $f(\infty)=q$, both interpreted as prime ends.
For the physical ensemble in $D$, let $W^f$ be its driving process in this
chart, with half-plane capacity normalization
\[
 \partial_t g_t(z)=\frac{2}{g_t(z)-W^f_t},\qquad W^f_0=0.
\]
Assume that it is a continuous semimartingale in its natural trace filtration,
with decomposition
\[
 W^f_t=M^f_t+\int_0^t b^f_s\,ds,\qquad
 C^f_t=\langle M^f\rangle_t.
 \tag{DR1}
\]
The drift-density paths $b^f$ are assumed finite and right-continuous,
after local stopping if necessary.  No absolute continuity or regularity
of $C^f$ is assumed.  At the initial time suppose
\[
 b^f_0=B(f)\in\mathbb R
 \tag{DR2}
\]
is deterministic.  It may depend on the entire physical marked domain and
its conformal chart.  In particular no Markov property of the single real
driver is assumed.  A finite deterministic initial variance rate $V(f)$
will be \emph{derived} below from locality and (DR1)--(DR2).

Let $A\subset\overline{\mathbb H}$ be a half-plane hull disjoint from $0$,
let $h:\mathbb H\setminus A\to\mathbb H$ be hydrodynamically normalized,
and set
\[
 \bar h(z)=h(z)-h(0),\qquad f_A=f\circ\bar h^{-1}.
 \tag{DR3}
\]
The source is still the same physical prime end, since $\bar h(0)=0$;
the target is also unchanged.  Put $d=h'(0)>0$, $e=h''(0)$.
Source locality, together with the usual deletion-chart calculation, will
give a finite deterministic initial variance rate $V(f)$ and the covariance
\[
 \boxed{\quad V(f_A)=V(f),\qquad
 B(f_A)=\frac{B(f)}{d}
       +\left(\frac{V(f)}2-3\right)\frac{e}{d^2}.\quad}
 \tag{DR4}
\]
These are identities for \emph{initial} characteristics and do not identify
the two complete driving laws.

Here are the details without assuming a quadratic-variation density.
If $h_t$ is the evolving deletion map, then $du=h_t'(W_t)^2dt$, and
It\^o's formula gives
\[
 d\widehat W_{u(t)}
  =h_t'(W_t)\,dM_t+h_t'(W_t)b_t\,dt
       +\frac12 h_t''(W_t)\,dC_t-3h_t''(W_t)\,dt.
 \tag{DR5}
\]
The $-3$ comes from
\[
 \partial_t h_t(z)
 =\frac{2h_t'(W_t)^2}{h_t(z)-h_t(W_t)}
       -\frac{2h_t'(z)}{z-W_t},\qquad
 (\partial_t h_t)(W_t)=-3h_t''(W_t).
\]
The comparison of filtrations is justified by geometric invertibility:
from a deleted-domain initial trace one reconstructs its original trace and
capacity, so the inverse capacity clock is adapted and its inverse is a
stopping-time clock.  Locality therefore compares the same trace filtration,
not an unrelated enlargement.  Adaptedness of an arbitrary increasing clock
would not suffice for this step.

\paragraph{Local bracket-regularization lemma (proved).}
Suppose $h''(0)\ne0$.  Then on a positive initial interval, $C_t$ is absolutely
continuous and has a finite right-continuous density.  In particular its
initial right derivative $V(f)$ exists and is deterministic.

\begin{proof}
The deleted-domain driver also has an absolutely continuous drift,
$\int\widehat b_u\,du$, by (DR1).  Uniqueness of the continuous
semimartingale decomposition applied to (DR5), after the justified capacity
time change, gives the signed-measure identity
\[
 h_t''(W_t)\bigl(dC_t-6\,dt\bigr)
   =2\bigl[h_t'(W_t)^2\widehat b_{u(t)}
                      -h_t'(W_t)b_t\bigr]dt.
 \tag{DR5a}
\]
The second derivative starts nonzero and remains so on a positive initial
interval.  On that interval the bracket is absolutely continuous, with
\[
 v_t=6+2\frac{h_t'(W_t)^2\widehat b_{u(t)}
                         -h_t'(W_t)b_t}{h_t''(W_t)}.
 \tag{DR5b}
\]
The derivative paths of $h_t$ along the continuous driving point are
continuous, while the two drift-density paths are right-continuous and
finite.  Thus $v$ is right-continuous and finite there.  In particular
\[
 V(f)=v_0
   =6+2\frac{d^2 B(f_A)-d B(f)}e
 \tag{DR5c}
\]
is deterministic.  This is an intrinsic initial rate of $C$, so it does
not depend on which nonzero-$e$ deletion was used to derive it.
Applying the same time change to bracket gives
$V(f_A)=V(f)$.  Rearranging (DR5c) gives the drift identity (DR4).
\end{proof}

There is always an available nonzero-$e$ test deletion, for example one
vertical slit at any fixed real point different from $0$.  Right-continuity
makes the initial drift and variance rates intrinsic, rather than arbitrary
versions of densities at a single time.  The bracket density has only been
proved locally at the start; its global regularity is not being presumed.

\paragraph{Source-neighborhood convention.}
In half-plane coordinates these slit deletions preserve a full neighborhood
of the source end.  Under $f$ this is a prime-end neighborhood.  Thus (DR4)
is justified on arbitrary rough domains if the corrected locality assumption
is formulated for such neighborhoods, or if these particular deletions are
explicitly admissible.  For regular analytic source boundaries it is ordinary
Euclidean neighborhood locality.  If locality is instead stated only using
Euclidean balls at non-singleton source impressions, the all-rough-domain
extension must remain an explicit additional hypothesis.

\subsection{Two vanishing slits force both characteristics}

For $\epsilon>0$ take the reflected vertical slits
\[
 A_\epsilon^\pm
  =\{\,\pm\epsilon^2+iy:0\leq y\leq2\epsilon\,\}.
 \tag{DR6}
\]
Their hydrodynamic maps are
\[
 h_\epsilon^\pm(z)
 =\pm\epsilon^2+
   \sqrt{(z\mp\epsilon^2)^2+4\epsilon^2},
 \tag{DR7}
\]
with the square-root branch asymptotic to $z\mp\epsilon^2$ at infinity.
Let $f_\epsilon^\pm=f\circ
(h_\epsilon^\pm-h_\epsilon^\pm(0))^{-1}$.
Both sequences are conformal charts of bounded simply connected domains and
\[
 f_\epsilon^\pm\longrightarrow f
 \quad\text{locally uniformly in }\mathbb H.
 \tag{DR8}
\]
Indeed the slits have diameter $2\epsilon$ and approach $0$; directly from
(DR7), their maps and inverse maps converge to the identity on compact
subsets of $\mathbb H$, while $h_\epsilon^\pm(0)\to0$.
Crucially their boundary derivatives do \emph{not} converge to the identity
jet:
\[
 d_\epsilon=(h_\epsilon^\pm)'(0)
     =\frac{\epsilon}{\sqrt{\epsilon^2+4}}\longrightarrow0,
 \qquad
 e_\epsilon^\pm=(h_\epsilon^\pm)''(0)
     =\mp\frac{4}{\epsilon(\epsilon^2+4)^{3/2}}.
 \tag{DR9}
\]
In particular
\[
 \frac{e_\epsilon^\pm}{d_\epsilon^2}
 =\mp\frac{4}{\epsilon^3\sqrt{\epsilon^2+4}}.
 \tag{DR10}
\]
Thus compact convergence of inverse charts permits very large second-jet
changes precisely at the source boundary.

\paragraph{Initial rigidity theorem (proved).}
Fix $f$ and assume (DR1)--(DR4) for $f$ and the charts $f_\epsilon^\pm$.
If
\[
 \sup_{0<\epsilon<\epsilon_0}
       \max\{|B(f_\epsilon^+)|,|B(f_\epsilon^-)|\}<\infty
 \tag{DR11}
\]
for some $\epsilon_0>0$, then
\[
 B(f)=0,\qquad V(f)=6.
 \tag{DR12}
\]
No continuity of $B$ is required.

\begin{proof}
Write $b=B(f)$, $v=V(f)$ and $b_\epsilon^\pm=B(f_\epsilon^\pm)$.
Equation (DR4) gives
\[
 b_\epsilon^\pm
 =\frac b{d_\epsilon}
    +\left(\frac v2-3\right)
                       \frac{e_\epsilon^\pm}{d_\epsilon^2}.
\]
Since $e_\epsilon^-=-e_\epsilon^+$,
\[
 b_\epsilon^++b_\epsilon^-=\frac{2b}{d_\epsilon},\qquad
 b_\epsilon^+-b_\epsilon^-=
       (v-6)\frac{e_\epsilon^+}{d_\epsilon^2}.
 \tag{DR13}
\]
The first right side is bounded only if $b=0$, since
$d_\epsilon\to0$.  The second is bounded only if $v=6$, since its multiplier
has magnitude asymptotic to $2\epsilon^{-3}$.  This proves (DR12).
\end{proof}

\paragraph{Exact instability and weaker rate conditions.}
For real numbers $U,C$ one has
$\max\{|U+C|,|U-C|\}=|U|+|C|$.  Hence the two drift covariances give the
exact identity
\[
 \boxed{\quad
 \max\{|b_\epsilon^+|,|b_\epsilon^-|\}
 =\frac{|b|\sqrt{\epsilon^2+4}}{\epsilon}
  +\frac{4|v/2-3|}{\epsilon^3\sqrt{\epsilon^2+4}}.
 \quad}
 \tag{DR14}
\]
There is no possible cancellation between the two reflected tests.
In particular, if the displayed maximum is $o(\epsilon^{-3})$, then $v=6$.
If it is $o(\epsilon^{-1})$, then $b=0$ and $v=6$.
Boundedness is more than sufficient.  Conversely, if $v\ne6$, its leading
blow-up is $|v-6|\epsilon^{-3}$; if $v=6$ and $b\ne0$, its leading blow-up
is $2|b|\epsilon^{-1}$.  This is a concrete, quantitative obstruction to any
alternative semimartingale driving law with tame initial drift under source
perturbations.

\subsection{Local boundedness in marked-chart topology suffices}

Equip inverse marked charts with locally uniform convergence on
$\mathbb H$, retaining the fixed preimages $0,\infty$ of source and target.
If $B$ is locally bounded at every $f$ in this topology, then (DR8) implies
(DR11).  In particular continuity of $B$ is sufficient, but unnecessarily
strong.  This is a hypothesis about the driving characteristic itself;
weak convergence of curve laws alone does not imply it.

\paragraph{Conditional classification theorem (proved).}
Suppose a Borel family of continuous Loewner-generating exploration laws on
all bounded simply connected marked prime-end domains has strong DMP and
source-neighborhood locality (including the test deletions (DR6) in each
chart).  Suppose its drivers satisfy (DR1)--(DR2), with finite right-continuous
drift-density paths and the usual restart compatibility of characteristic
kernels, with no initial assumption on quadratic variation.  If the initial drift map $B$ is locally bounded in marked inverse-
chart topology, then every driving process is
\[
 W^f_t=\sqrt6\,B_t,
\]
and the family is chordal SLE$_6$ in these conformal coordinates.

\begin{proof}
The initial rigidity theorem gives $B(f)=0$ and the derived $V(f)=6$ for
every marked inverse chart.  At any deterministic rational time $r$,
condition on the initial exploration and use strong DMP in its remaining
marked domain and current conformal chart.  The centered future driver is
exactly the restarted driver; its initial drift is $b^f_r$, by
right-continuity.  The restarted law has initial drift zero, so
$b^f_r=0$ almost surely.  Countably many rational times and right-continuity
give $b^f_t=0$ throughout each localized lifetime.  Therefore all domain
chart drivers are continuous local martingales.

For completeness the no-density local-martingale locality argument then
identifies their brackets.  With both drifts zero, (DR5a) becomes
$h_t''(W_t)(dC_t-6\,dt)=0$.  A nonzero-$e$ vertical deletion yields
$C_t=6t$ on a positive initial interval.  Let
\[
 T=\inf\{t\geq0:C_t-6t\ne0\}.
\]
This is an entrance stopping time of a continuous adapted process into an
open set, in the usual right-continuous filtration.  The initial zero
interval gives $T>0$.  If $T<\infty$, continuity gives $C_T=6T$ and the
first-deviation definition gives nonzero values arbitrarily soon after $T$.
Strong DMP restarted at $T$ gives another positive zero interval, a
contradiction.  Thus $C_t=6t$ throughout.  L\'evy's characterization gives
Brownian driving, and Loewner generation identifies SLE$_6$.
\end{proof}

The Borel-kernel condition means that the family admits the standard strong
conditional-kernel interpretation on random slit domains, and that the
characteristics can be chosen consistently on that family.  It is not a
claim that deterministic-domain convergence of discrete explorations proves
this property.  Translation invariance is not needed in this conditional
theorem.  If its hypotheses hold only after one fixed physical affine
normalization, the conclusion is the corresponding fixed-affine SLE$_6$
family; an affine change of physical coordinates is not presumed to preserve
semimartingale driving in an arbitrary preselected complex structure.

\subsection{A finite-dimensional closure that is completely classified}

For example, fix finitely many interior points $z_j\in\mathbb H$ and
finitely many nonnegative derivative orders $m_j$.  Let
\[
 J(f)=\bigl(f^{(k)}(z_j):1\leq j\leq N,
                              0\leq k\leq m_j\bigr).
\]
Suppose the initial drift has a representation
\[
 B(f)=\mathcal B(J(f)),
 \tag{DR15}
\]
where $\mathcal B$ is locally bounded at every realized jet vector.
The bracket has no assumed density or finite-dimensional representation;
its initial rate is supplied by the local regularization lemma.
Locally uniform convergence of conformal maps implies convergence of all
these interior jets by Cauchy's formula, so (DR15) implies local boundedness
of $B$ in marked-chart topology.  The conditional classification theorem
therefore applies: within (DR1)--(DR2), every source-local strong-DMP family
with such a locally bounded finite-dimensional \emph{bulk-jet drift closure}
is SLE$_6$.

This statement does not cover a closure depending on unbounded source
boundary derivatives such as $f'(0)$ or a boundary pre-Schwarzian.  Those
quantities can change without control under (DR8), and can carry exactly the
singular terms in (DR13).  Nor does it show that true percolation limits have
semimartingale driving in a fixed physical complex structure.  The new
rigorous reduction is that any non-SLE candidate with regular characteristic
rates must violate the displayed near-source drift bound, in the explicit
$\epsilon^{-1}$ or $\epsilon^{-3}$ manner above.


\subsection{The criterion also detects a wrong affine coordinate system}

\paragraph{Corollary (proved).}
Let $A$ be an orientation-preserving real-linear map which is not a Euclidean
similarity, and let the physical ensemble be the fixed-affine pullback
\[
 P_{D,p,q}=(A^{-1})_*
       \operatorname{Law}(\mathrm{SLE}_6(AD;Ap,Aq)).
\]
Although this family satisfies corrected source-neighborhood locality and
strong DMP, in the original Euclidean complex structure it cannot satisfy
both the regular semimartingale characteristic assumptions (DR1)--(DR2) in
all marked domains and the locally bounded initial-drift condition above.

\begin{proof}
If both properties held, the conditional classification theorem would make
this ensemble ordinary SLE$_6$ in the original complex structure.  The
single-domain affine identifiability theorem would then make $A$ a
similarity, a contradiction.
\end{proof}

Thus the drift criterion selects a conformal structure.  It is not a
coordinate-free consequence of coarse crossing bounds: bounded affine
changes preserve such bounds, whereas a wrong affine structure necessarily
violates at least one of the criterion's additional analytic hypotheses.
The corollary does not decide whether the failure is semimartingality,
right-continuity of the drift density, deterministic initial drift,
or boundedness of the initial drift; distinguishing these failures is a
separate question.

\subsection{What corrected locality says at a regular analytic source germ}

Here is a separate necessary local structure statement that uses the full
corrected locality, including changes of the distant target.  It concerns
initial characteristics only, not existence of a process from arbitrary
coefficient choices.

Suppose the physical boundary is analytic near $p$ and the source colors
are fixed there.  Let $G$ denote this embedded colored source germ.  For a
chart $f$ extending conformally across $0$, define
\[
 a(f)=|f'(0)|>0,\qquad
 P(f)=\operatorname{Re}\frac{f''(0)}{f'(0)}.
 \tag{DR16}
\]
If another physical marked domain has the same source germ but any distant
target, let $k$ be its inverse chordal chart with $k(0)=p$.  The local map
$h=k^{-1}\circ f$ is analytic across $0$, takes the real boundary interval
to a real interval, and has $h(0)=0$, $d=h'(0)>0$, $e=h''(0)\in\mathbb R$.
The two initial physical traces agree by corrected locality.  The local
Loewner coordinate comparison is exactly (DR4), with the same derivation
from pole cancellation and It\^o's formula; a globally defined deletion map
is unnecessary.  Consequently $V(f)=V(k)$ depends only on $G$.

\paragraph{Proposition (proved).}
Under the semimartingale assumptions above, corrected locality at analytic
source germs is equivalent at the level of initial characteristic covariance
to the existence of two source-germ quantities $V(G)$ and $\beta(G)$ such that
\[
 \boxed{\quad
 V(f)=V(G),\qquad
 B(f)=a(f)\,\beta(G)
                 +\left(3-\frac{V(G)}2\right)P(f).
 \quad}
 \tag{DR17}
\]
Here ``equivalent at the level of covariance'' means precisely that every
covariant initial-characteristic assignment has this form, and every
assignment of this form obeys the local chart covariance identity.  It does
not mean that arbitrary $V,\beta$ yield a domain Markov exploration family.

\begin{proof}
Since $k=f\circ h^{-1}$, the chain rule gives
\[
 a(k)=\frac{a(f)}d,\qquad
 P(k)=\frac{P(f)}d-\frac e{d^2}.
\]
Write $c=V(G)/2-3$.  Combining these transformations with
$B(k)=B(f)/d+c e/d^2$ yields
\[
 \frac{B(k)+cP(k)}{a(k)}
       =\frac{B(f)+cP(f)}{a(f)}.
\]
This common value is the source-germ invariant $\beta(G)$, proving (DR17).
Conversely direct substitution gives the required covariance.
\end{proof}

Physical translation invariance makes $V(G)$ and $\beta(G)$ invariant under
translations of the colored germ.  If dilation covariance is separately
assumed, then $V(\lambda G)=V(G)$ and
$\beta(\lambda G)=\lambda^{-1}\beta(G)$.  For a straight source-boundary germ,
which is unchanged up to translation by dilation, this would force
$\beta(G)=0$; the dimensionless variance can still depend on the boundary
orientation before an affine metric has been identified.

The boundary pre-Schwarzian term in (DR17) explains the quantitative anomaly
in the vanishing-slit test.  Even when the physical source germ is unchanged,
$P(f_\epsilon)$ may grow as $\epsilon^{-3}$ while $f_\epsilon\to f$ on every
interior compact set.  Thus source locality by itself gives a connection-type
transformation law, rather than vanishing drift.  The locally bounded bulk
closure theorem removes exactly this uncontrolled boundary dependence.
The missing global step is to understand how DMP constrains these local
assignments at rough, dynamically created source germs, where the analytic
boundary jet in (DR16) no longer exists.

% Research fragment; requires amsmath/amssymb/amsthm.
\section{A wrong affine metric cannot have tame driving drift in every bounded domain}
\label{sec:scaling-semimartingale}

\paragraph{Status.}
The conclusion here is stronger than the earlier wrong-coordinate consequence
of drift boundedness.  For a fixed non-similarity affine deformation of SLE$_6$,
regular semimartingale driving with absolutely continuous right-continuous
drift cannot hold in every bounded domain in the wrong Euclidean complex
structure.  No local boundedness of the initial drift functional is needed.
The proof uses exact scale invariance only for the \emph{known fixed-affine
SLE$_6$ half-plane law}; it does not assume that arbitrary critical
subsequential limits are scale invariant.  All characteristic regularity
assumptions remain explicit.

\subsection{Exact Brownian scaling and a regular initial bracket}

\paragraph{Scaling lemma (proved).}
Let $W$ be a continuous real process on $[0,\infty)$ with $W_0=0$, whose law
has exact Brownian scaling:
\[
 (c^{-1/2}W_{ct})_{t\geq0}\ \stackrel{\rm law}{=}\ (W_t)_{t\geq0}
 \qquad(c>0).
 \tag{SM1}
\]
Suppose up to a strictly positive stopping time it is a continuous
semimartingale
\[
 W_t=M_t+\int_0^t b_s\,ds,
 \tag{SM2}
\]
where $M_0=0$ is a continuous local martingale and $b$ is finite and
right-continuous at $0$.  Suppose its bracket has a finite deterministic
initial right derivative
\[
 \lim_{t\downarrow0}\frac{\langle M\rangle_t}{t}=v\in[0,\infty)
 \quad\text{almost surely}.
 \tag{SM3}
\]
Then the entire law of $W$ is that of $\sqrt v B$.  This includes the
identically zero process if $v=0$.

\begin{proof}
Fix $T<\infty$ and stop the local decomposition at a positive initial
localizing time.  As $c\downarrow0$, with probability tending to one the
whole interval $[0,cT]$ lies before that time.  Right-continuity of $b$ gives
\[
 \sup_{t\leq T}\left|c^{-1/2}\int_0^{ct}b_s\,ds\right|
 \leq\sqrt c\,T\sup_{s\leq cT}|b_s|\longrightarrow0
\]
almost surely on the localized event.  Writing $C_t=\langle M\rangle_t$,
(SM3) implies
\[
 \sup_{t\leq T}|C_{ct}/c-vt|\longrightarrow0.
 \tag{SM4}
\]
Indeed $|C_{ct}/c-vt|\leq t\sup_{0<s\leq cT}|C_s/s-v|$.
The continuous-martingale functional central limit theorem now gives
$c^{-1/2}M_{c\cdot}\Rightarrow\sqrt v B$ on $C[0,T]$.
One may see this directly from Dambis--Dubins--Schwarz: extend its Brownian
motion if necessary and write
\[
 c^{-1/2}M_{ct}=B^{(c)}_{C_{ct}/c},
 \qquad B^{(c)}_u=c^{-1/2}\widetilde B_{cu}.
\]
Each $B^{(c)}$ has standard Brownian law.  Uniform bracket convergence to the
deterministic clock and the Brownian modulus of continuity imply the stated
limit; independence between $B^{(c)}$ and the bracket is unnecessary.
Adding the negligible drift gives the same limit for $W$.
By exact scaling, however, each rescaled process already has the law of
$W$ on $[0,T]$.  That law is therefore Brownian with variance $v$.
Arbitrariness of $T$ proves the assertion.
\end{proof}

The regular initial bracket is not automatic for a general self-similar
continuous process.  In the affine-SLE application below it follows from
locality and absolutely continuous drift in two domain charts, using the
local bracket-regularization identity.

\subsection{The half-plane affine-SLE obstruction}

Let $A$ be an invertible orientation-preserving real-linear map and let
$\eta$ be chordal SLE$_6$ in $A\mathbb H$ from $0$ to $\infty$.  Put
$\gamma=A^{-1}\eta\subset\overline{\mathbb H}$ and let $W$ be its
half-plane-capacity driving function.  Its continuous Loewner generation and
infinite capacity lifetime are supplied by the affine-admissibility result
already proved in the research record.  Its law has exact Brownian scaling
(SM1): ordinary SLE is dilation covariant, $A$ commutes with positive scalar
dilations, and capacity scales by the square of physical length.

\paragraph{Proposition (proved).}
If this $W$ satisfies (SM2)--(SM3), then $A$ is a Euclidean similarity.

\begin{proof}
The scaling lemma makes $W$ ordinary Brownian driving of variance $v$.
Thus the trace is ordinary chordal SLE$_v$ in $\mathbb H$.
We identify $v$ without a Hausdorff-dimension theorem.

For $\kappa\leq4$ the nonzero real boundary points are never swallowed by
SLE$_\kappa$: their Loewner distances divided by $\sqrt\kappa$ are Bessel
processes of dimension $1+4/\kappa\geq2$ (the case $\kappa=0$ is immediate).
An affine image of SLE$_6$ swallows boundary points, so $v>4$.
For any $\kappa>4$, the boundary-race calculation gives
\[
 \mathbb P_\kappa(\tau_y<\tau_x)
 =F_\kappa\!\left(\frac{-x}{y-x}\right),\qquad
 F_\kappa(r)=
 \frac{\displaystyle\int_0^r[u(1-u)]^{-4/\kappa}\,du}
 {\displaystyle\int_0^1[u(1-u)]^{-4/\kappa}\,du},
 \quad x<0<y.
 \tag{SM5}
\]
This includes $\kappa\geq8$.  To check it, the ratio
$R=-X/(Y-X)$ has generator
\[
 \frac\kappa2\frac{d^2}{dr^2}
   +\frac{2(1-2r)}{r(1-r)}\frac d{dr}
\]
after $ds=dt/(Y-X)^2$.  Its scale density is
$[r(1-r)]^{-4/\kappa}$, integrable at both ends exactly when $\kappa>4$.
The first swallowing is finite almost surely by the Bessel description;
simultaneous swallowing is impossible since $Y-X$ is positive and
nondecreasing.  Bounded optional stopping gives (SM5).

Rotate $A\mathbb H$ conformally back to $\mathbb H$.  The resulting map
has matrix
\[
 T=\begin{pmatrix}\alpha&\beta\\0&\delta\end{pmatrix},
 \qquad \alpha,\delta>0.
\]
It acts on the real boundary by $x\mapsto\alpha x$, so it preserves the
ratio $-x/(y-x)$ and the topological swallowing order.  The law of
$\gamma$ therefore has boundary-race function $F_6$, while its Brownian
driving gives $F_v$.  Equality for every $r\in(0,1)$ forces $v=6$:
compare the derivatives in (SM5), proportional respectively to
$[r(1-r)]^{-4/v}$ and $[r(1-r)]^{-2/3}$.

Finally $T^{-1}\mathrm{SLE}_6(\mathbb H;0,\infty)$ is ordinary SLE$_6$.
Half-plane affine identifiability now implies $T$ is a similarity.
For completeness, the strictly increasing SLE$_6$ left-passage function
\[
 L(u)=\frac{\int_{-\infty}^u(1+s^2)^{-2/3}\,ds}
 {\int_{-\infty}^{\infty}(1+s^2)^{-2/3}\,ds}
\]
and preservation of topological side give
$L(u)=L((\alpha u+\beta)/\delta)$ for all real $u$.
Thus $\alpha=\delta$ and $\beta=0$.  Rotation does not affect whether $A$
is a similarity, proving the proposition.
\end{proof}

\subsection{Transfer from bounded-domain regularity to the half-plane germ}

The preceding proposition is already a complete unbounded-domain statement.
The following additional argument shows why requiring the tame drift
structure in \emph{every bounded domain} is enough to trigger it.
We spell out the local coordinate transfer, rather than silently assume
semimartingality is preserved by a general physical affine map.

\paragraph{Bounded-domain obstruction theorem (proved).}
Consider the canonical fixed-affine SLE$_6$ family
\[
 P_{D,p,q}=(A^{-1})_*\operatorname{Law}
       (\mathrm{SLE}_6(AD;Ap,Aq))
\]
on bounded simply connected marked domains.  Suppose in the original
Euclidean complex structure every domain-chart driver is a continuous
semimartingale whose finite-variation part is
$\int_0^t b_s\,ds$, with finite right-continuous drift density and
deterministic initial value $B(f)$, consistently under the usual conditional
kernel interpretation.  Then $A$ is a Euclidean similarity.
No boundedness or continuity of the map $f\mapsto B(f)$ is assumed.

\begin{proof}
Choose a bounded half-disk $D\subset\mathbb H$ with $0$ in the interior of
its real diameter, and a distinct target on that diameter away from $0$.
Choose the boundary colors so they agree near $0$ with the half-plane
problem $(\mathbb H,0,\infty)$.  Let $\phi:D\to\mathbb H$ be its chordal
chart sending these two marks to $0,\infty$.  It extends analytically
across $0$, with $\phi'(0)>0$.

The local bracket-regularization lemma applies to this bounded-domain
driver: delete one further vertical slit in its $\phi$ chart, away from its
marks, with nonzero $h''(0)$.  Both domains have tame absolutely continuous
drift by hypothesis.  The signed-measure identity
\[
 h_t''(W_t)(dC_t-6\,dt)
  =2\bigl[h_t'(W_t)^2\widehat b_{u(t)}
                        -h_t'(W_t)b_t\bigr]dt
\]
therefore gives a finite right-continuous initial bracket density in $D$,
with deterministic initial value $v_D$.

Corrected source-neighborhood locality of ordinary SLE$_6$, preserved by
the fixed map $A^{-1}$, couples this bounded-domain curve to the canonical
affine-SLE$_6$ half-plane curve so that their physical initial segments are
identical until leaving a sufficiently small half-neighborhood of $0$.
This uses ordinary SLE$_6$ locality including distant target changes; it is
not a new hypothesis about an unknown nonconformal family.
On that common initial trace, the coordinate change is the analytic map
$\phi$ at the source.

A local analytic coordinate change and its inverse preserve the tame
semimartingale structure just obtained.  More explicitly, if $t$ is physical
half-plane capacity and $u$ is the bounded-chart capacity, the evolving local
map satisfies
\[
 W_D(u(t))=h_t(W_H(t)),\qquad
 \dot u(t)=h_t'(W_H(t))^2.
\]
Initially $h_0=\phi$, so $h_0'(0)>0$ and all required local derivatives are
finite.  Before the trace leaves the common neighborhood, localize so the
first derivative stays bounded away from zero and the derivatives stay
bounded.  Both capacity clocks are adapted functions of the same initial
trace; their inverses are stopping-time clocks.  The inverse field
$h_t^{-1}$ has finite variation in time, as follows by differentiating the
Loewner composition.  Applying the finite-variation It\^o--Wentzell formula
to the semimartingale $W_D$, and changing the capacity clock, produces a
continuous-semimartingale decomposition of $W_H$ on a positive initial
interval.  Its drift is absolutely continuous and right-continuous, and its
initial bracket density is the same scalar $v_D$.  These conclusions use
only a \emph{conformal local chart transition}, not the nonconformal map $A$.

For an explicit initial check, writing $d=\phi'(0)$, $e=\phi''(0)$,
the transferred rates obey
\[
 V_H=v_D,\qquad
 B_H=dB_D-\left(\frac{v_D}{2}-3\right)\frac ed,
 \tag{SM6}
\]
so both are finite and deterministic.  Thus the canonical half-plane driver
satisfies (SM2)--(SM3).  Its exact scaling and the half-plane proposition
force $A$ to be a similarity.
\end{proof}

Only a bounded domain flat at its source and one of its source-preserving
slit deletions are needed for the bracket part of this contradiction.
The full all-bounded-domain assumption is a convenient sufficient
formulation.  The obstruction does not identify the precise failure for a
particular wrong affine metric: semimartingality, absolute continuity or
right-continuity of the drift, or deterministic initial drift may fail.
It does show that at least one must fail somewhere.  It supplies no argument
that an unknown exploration family admits a correct affine normalization
with tame drift.

\section{An affine metric is identifiable from one chordal law}
\label{sec:identifiability}

For a bounded Jordan domain $D$, distinct $p,q\in\partial D$, and an
orientation-preserving invertible real-linear map $A$, put
\begin{equation}
 \nuA{A}{D}{p}{q}
 = (A^{-1})_*\Law\bigl(\SLE(AD;Ap,Aq)\bigr).
 \label{eq:affine-law}
\end{equation}
Translations in an affine normalizer are immaterial. The tensor
$A^{\mathsf T}A$ is its Euclidean conformal metric, determined only up to a
positive scalar.

\subsection{A boundary-swallowing race}

\begin{lemma}[Strict boundary-race probability]
\label{lem:race}
Let $\eta$ be chordal SLE$_6$ in $\HH$ from $0$ to $\infty$. For
$x<0<y$, let $\tau_x,\tau_y$ denote their Loewner swallowing times. Then
\begin{equation}
 \mathbb P(\tau_y<\tau_x)=F\left(\frac{-x}{y-x}\right),\qquad
 F(r)=\frac{\int_0^r[u(1-u)]^{-2/3}\,du}
 {\int_0^1[u(1-u)]^{-2/3}\,du}.
 \label{eq:race}
\end{equation}
In particular $F$ is strictly increasing on $(0,1)$. The event in
\eqref{eq:race} is invariant under orientation-preserving homeomorphisms of
marked Jordan domains, when formulated as the order of separation from
the terminal prime end by the filled exploration hull.
\end{lemma}

\begin{proof}
Use $\partial_tg_t(z)=2/(g_t(z)-W_t)$ with $W_t=\sqrt6 B_t$. Until the
first swallowing time put
\[
 X_t=g_t(x)-W_t<0,\quad Y_t=g_t(y)-W_t>0,\quad
 T_t=Y_t-X_t,\quad R_t=-X_t/T_t\in(0,1).
\]
Then
\begin{align*}
 dX_t&=\frac{2}{X_t}\,dt-\sqrt6\,dB_t,&
 dY_t&=\frac{2}{Y_t}\,dt-\sqrt6\,dB_t,\\
 dT_t&=\frac{2T_t}{-X_tY_t}\,dt>0,&
 dR_t&=\frac{2(1-2R_t)}{R_t(1-R_t)T_t^2}\,dt
                 +\frac{\sqrt6}{T_t}\,dB_t.
\end{align*}
With $ds=dt/T_t^2$, the generator of $R$ is
\[
 \mathcal L=3\frac{d^2}{dr^2}
       +\frac{2(1-2r)}{r(1-r)}\frac d{dr}.
\]
The increasing solution of $\mathcal LF=0$, $F(0)=0$, $F(1)=1$ has
$F'(r)=C[r(1-r)]^{-2/3}$, giving \eqref{eq:race} by bounded optional
stopping. To justify the terminal values, $|X|/\sqrt6$ and $Y/\sqrt6$
are Bessel processes of dimension $1+4/6=5/3$ up to their first zero;
thus the first swallowing time is finite almost surely. Alternatively,
the accessible endpoints of the displayed diffusion give the same stopping
argument. There is no simultaneous swallowing, since $T_t\geq y-x>0$,
whereas both $X_t,Y_t\to0$ would force $T_t\to0$. At the first swallowing,
$R=1$ exactly when $y$ is swallowed first. ODE solutions and the continuous
driver stay finite up to that first time.

Swallowing means that the marked boundary end has been cut off from the
target by the explored filled hull (including the boundary swallowed at
first contact). This definition uses only the domain, the oriented path,
and its order of traversal. A homeomorphism carries filled target
components to filled target components. Capacity parametrization therefore
has no role in the comparison event.
\end{proof}

\subsection{Boundary rigidity and affine uniqueness}

\begin{lemma}[A single SLE law identifies its conformal boundary map]
\label{lem:boundary-rigidity}
Let $h:\overline D\to\overline{D'}$ be an orientation-preserving
homeomorphism between closed bounded Jordan domains, with
$p'=h(p)$, $q'=h(q)$. If $h_*\SLE(D;p,q)=\SLE(D';p',q')$, then
$h|_{\partial D}$ agrees with the boundary extension of a conformal
map $f:D\to D'$.
\end{lemma}

\begin{proof}
Choose conformal maps $\phi:D\to\HH$ and $\phi':D'\to\HH$ sending
their marked endpoints to $0,\infty$. Their Jordan boundary extensions
give an increasing homeomorphism
\[
 H=\phi'\circ h\circ\phi^{-1}:\overline\RR\to\overline\RR
\]
fixing $0,\infty$. For every $x<0<y$, equality of the corresponding
boundary-race probabilities and strictness in Lemma~\ref{lem:race} give
\[
 \frac{-H(x)}{H(y)-H(x)}=\frac{-x}{y-x},\quad\text{hence}\quad
 \frac{H(x)}{H(y)}=\frac{x}{y}.
\]
Fixing one $x<0$ and one $y>0$ shows $H(s)=cs$ for all real $s$, with
$c>0$. Consequently $f=(\phi')^{-1}\circ(c\phi)$ has exactly the
required boundary values. Equality of interior maps is not asserted for
an arbitrary $h$.
\end{proof}

\begin{theorem}[Single-domain affine identifiability]
\label{thm:affine-identifiability}
For even one bounded Jordan $D$ and one pair of distinct boundary points
$p,q$, the equality $\nuA{A}{D}{p}{q}=\nuA{B}{D}{p}{q}$ holds if and
only if $BA^{-1}$ is a Euclidean similarity. Equivalently,
\[
 B^{\mathsf T}B=c A^{\mathsf T}A\qquad(c>0).
\]
The statement also allows orientation-reversing normalizers by composing
one map with a reflection.
\end{theorem}

\begin{proof}
The equality of laws means $M=BA^{-1}$ carries the SLE law in $AD$ to
the SLE law in $BD$, at the corresponding endpoints. By
Lemma~\ref{lem:boundary-rigidity}, a conformal map $f:AD\to BD$ has
boundary values equal to $M$. Both real coordinate functions of the
real-linear map $M$ are harmonic. The real and imaginary parts of $f$
are also harmonic, continuous on the closure, with the same boundary
values. The maximum principle on the bounded domain gives $f=M$
throughout $AD$. Thus $M$ is complex-linear and invertible, hence a
similarity. The converse is ordinary SLE similarity covariance.
Reflection covariance reduces the remaining orientation case to this one.
\end{proof}


\begin{thebibliography}{9}
\bibitem{lawler}G. F. Lawler, \emph{Schramm--Loewner Evolution},
Park City lecture notes (2007), Lectures 3--4.
\url{https://math.uchicago.edu/~lawler/utah.pdf}.
\bibitem{yuan}Y. Yuan, \emph{Topological characterisations of Loewner traces},
Indiana Univ. Math. J. \textbf{71} (2022), 1027--1046.
\url{https://arxiv.org/abs/2003.05535}.
\bibitem{record}\emph{Research record: locality, terminal information, and planar exploration laws},
\texttt{progress.tex} in this repository, including the original-locality
investigation and the corrected-locality continuation (26 October 2026).
\end{thebibliography}
\end{document}
